\documentclass[12pt,a4paper]{article}
\usepackage[utf8]{inputenc}
\usepackage[T1]{fontenc}
\usepackage[hidelinks]{hyperref}
\usepackage{geometry}
\geometry{top=2.5cm, bottom=2.5cm, left=2.5cm, right=2.5cm}
\usepackage{titlesec}
\usepackage{enumitem}

\titleformat{\section}{\large\bfseries}{\thesection.}{1em}{}

\begin{document}

\begin{center}
    {\huge \textbf{TÈM DE REFERANS (TDR)}} \\[0.5cm]
    {\large \textbf{Pòs: Data Analyst \& Software Developer}} \\[0.2cm]
    {\small \textit{Direksyon Teknoloji ak Jesyon Done}}
\end{center}

\vspace{1cm}

\section{KONTEKS PWOJÈ A}
Òganizasyon an itilize yon sistèm koleksyon done ki baze sou platfòm ONA (\url{Ona.io}). Done sa yo alimente yon baz done MySQL ki deplwaye sou Hostinger. Yon aplikasyon ki devlope nan langaj PHP (pure PHP) responsab pou resevwa done yo atravè yon Webhook epi mete yo nan tablo respektiv yo. Pou asire vizibilite ak transparans, sistèm nan itilize scripts AWS pou sinkwonize done yo chak inèdtan sou Google Spreadsheets (Google Drive), epi Looker Studio sèvi kòm platfòm prensipal pou prezante Dashboards aktivite yo.

\section{OBJEKTIF TRAVAY LA}
Objektif prensipal misyon sa a se asire antretyen, amelyorasyon ak evolisyon sistèm koleksyon ak analiz done a. Kandida a dwe asire sistèm nan ap fonksyone san pwoblèm toutan, pandan l ap prepare tranzisyon vè yon achitekti ki pi solid ak modèn, si sa parèt nesesè.

\section{RESPONSABLITE PRENSIPAL}

\subsection{Devlopman ak Antretyen Lojisyèl}
\begin{itemize}
    \item Asire bon fonksyònman Webhook ki konekte ONA ak aplikasyon APA a.
    \item Rezoud tout bug oswa pwoblèm teknik ki ka parèt nan aplikasyon APA a ki deplwaye sou Hostinger la.
    \item Jere ak optimize baz done MySQL la pou pèfòmans maksimòm.
\end{itemize}

\subsection{Analiz Done (Data Analyst)}
\begin{itemize}
    \item Prevansyon erè: Siveye antre done yo pou evite duplikasyon oswa erè fòma.
    \item Kontwòl Kalite (Data Quality): Validate done yo pou asire entegrite enfòmasyon yo.
    \item Analiz Esploratwa (EDA): Idantifye tandans nan done yo pou ede nan pran desizyon.
\end{itemize}

\subsection{Otomatizasyon ak Rapò}
\begin{itemize}
    \item Siveye ak ajiste scripts AWS yo pou asire Google Spreadsheets yo ajou chak inèdtan.
    \item Kenbe ak amelyore Dashboards Looker Studio yo pou bay yon vizyon klè sou tout aktivite yo.
\end{itemize}

\section{ANTRETYEN, AMELYORASYON AK MODÈNIZASYON}
Pou asire sistèm nan toujou pèfòman epi fasil pou jere nan tan k ap vini yo, **Data Analyst / Software Developer** la gen responsablite ak libète konplè pou l chwazi pi bon metodoloji pou l amelyore ak modènize aplikasyon APA a. Sa gen ladan l :

\begin{itemize}
    \item \textbf{Entèvni ak Debogaj:} Responsablite pou l entèvni rapidman epi rezoud tout pwoblèm (bugs) ki parèt nan sistèm nan pandan tout dire kontra a.
    \item \textbf{Modènizasyon Achitekti:} Libète pou l kòmanse ak reyalize konvèsyon kòd la vè yon achitekti modèn (tankou Python FastAPI oswa lòt) pou pi bon pèfòmans ak sekirite (Scalability).
    \item \textbf{Dokimantasyon:} Obligasyon pou l dokimante nèt tout fonksyònman sistèm nan, estrikti kòd la, ak tout chanjman oswa amelyorasyon ki fèt pou asire transparans ak perennite pwojè a.
    \item \textbf{Amelyorasyon Kontini:} Kapasite pou l pwopoze epi aplike nouvo fonksyonalite ki ka rann travay analiz done a pi efikas.
\end{itemize}


\section{PROFIL KANDIDA A}
\begin{itemize}
    \item Eksperyans nan devlopman Web ak PHP ak Python (FastAPI).
    \item Metriz SQL (MySQL) ak jesyon baz done sou Hostinger.
    \item Konesans nan otomatizasyon ak scripts (AWS, API Google).
    \item Kapasite nan analiz done ak kreyasyon vizyalizasyon (Looker Studio).
    \item Lespri inisyativ pou pwopoze ak egzekite chanjman estriktirèl nesesè.
\end{itemize}

\vspace{1.5cm}

\begin{flushleft}
    Siyen nan Pòtoprens, Ayiti. \\
    Dat: \today
\end{flushleft}

\end{document}
